\documentclass[a4paper,14pt]{extarticle}
\usepackage[T2A]{fontenc}
\usepackage[utf8]{inputenc}
\usepackage[russian]{babel}
\usepackage{tempora,setspace,indentfirst}
\usepackage[left=30mm,right=15mm,top=20mm,bottom=20mm]{geometry}
\onehalfspacing
\setlength{\parindent}{1.25cm}
\begin{document}
\noindent\textit{Учебный пример. Условия задачи и данные вымышлены.}
\tableofcontents
\clearpage
\section*{\centering\normalsize ВВЕДЕНИЕ}
\addcontentsline{toc}{section}{Введение}
В учебной модели подразделения сведения о заявках распределены между отдельными
записями. Это затрудняет поиск обращения и определение его текущего состояния.
Актуальность темы в заданных условиях связана с потребностью представить
регистрацию, изменение статуса и поиск заявок в едином программном интерфейсе.
Данное описание является вымышленным условием задачи, а не результатом
обследования реальной организации.

Объект работы --- процесс учета заявок в условном подразделении. Предмет работы
--- модель данных и программные средства регистрации, изменения статуса
и поиска заявок в пределах учебного прототипа.

Цель работы --- разработать и функционально проверить прототип учета заявок.
Для достижения цели поставлены следующие задачи:
\begin{enumerate}
    \item определить требования к регистрации, изменению статуса и поиску заявок;
    \item спроектировать модель заявки с тремя статусами и допустимыми переходами;
    \item разработать прототип, реализующий указанные группы функций;
    \item проверить функции на учебном стенде по десяти сценариям.
\end{enumerate}

В работе предполагается использовать анализ сценариев, моделирование данных,
прототипирование и функциональное тестирование. Исходными материалами служат
условное описание процесса и искусственно созданные записи заявок, не содержащие
данных реальной организации. Для проверки предусмотрены четыре сценария
регистрации, три сценария изменения статуса и три сценария поиска.
Для каждого сценария задаются входные данные и ожидаемый результат.

Проверка ограничена одним пользователем и учебными данными. Нагрузочные испытания,
разграничение прав доступа, внедрение и оценка сокращения времени обработки
заявок не входят в границы работы. Предполагаемая практическая значимость состоит
в возможности использовать прототип как основу для последующей опытной эксплуатации;
ее результативность в рамках этой работы не подтверждается.

Основная часть включает три главы: анализ задачи и требований, проектирование
и разработку прототипа, функциональную проверку. В заключении планируется
сопоставить полученные результаты с поставленными задачами, указать ограничения
и направления дальнейшей работы.
\end{document}
